\documentclass[a4paper,12pt]{article}
\usepackage[utf8]{inputenc}
\usepackage[T1]{fontenc}
\usepackage[ngerman]{babel}
\usepackage{amsmath}
\usepackage{amssymb}
\usepackage{enumitem}
\usepackage{geometry}
\geometry{a4paper, margin=2.5cm}
\usepackage{parskip}

\title{Einsendeaufgaben -- Aufgabe 3.4\\[0.5em]
\large Modul 61112: Lineare Algebra}
\author{}
\date{}

\begin{document}
\maketitle

\section*{Aufgabe 3.4}

\begin{enumerate}[label=\alph*)]

    \item
    $C^2$ ist symmetrisch, da
    \[
      C^2 = -C(-C) = C^T C^T = (CC)^T = (C^2)^T.
    \]
    Nach Lemma 3.2.6 ist demnach $\beta$ symmetrisch.

    \item
    Es gilt:
    \begin{align*}
      \beta(v,v) &= v^T C^2 v \\
                 &= (v^T C)(C v) \\
                 &= (C^T v)^T (C v) \\
                 &= (-C v)^T (C v) \\
                 &= -(C v)^T (C v) \\
                 &\leq 0 \quad (\text{da } v^T v \geq 0\ \forall v \in \mathbb{R}^n)
    \end{align*}

    Wir untersuchen nun eine beliebige diagonale Matrixdarstellung
    $M_B(\beta) = \operatorname{diag}(\beta(v_1, v_1), \ldots, \beta(v_n, v_n))$
    mit $B = (v_1, \ldots, v_n)$. Für jedes $v_i$ mit $1 \leq i \leq n$ muss
    aus den vorherigen Überlegungen entweder $\beta(v_i, v_i) = 0$ oder
    $\beta(v_i, v_i) < 0$ gelten. Wir betrachten die Unterräume mit den Basen
    $B_1 := \{v_i \in B \mid \beta(v_i, v_i) = 0 \text{ für alle } 1 \leq i \leq n\}$
    und $B_2 := \{v_i \in B \mid \beta(v_i, v_i) < 0 \text{ für alle } 1 \leq i \leq n\}$.

    Wir definieren $f: V \to V,\ v \mapsto C v$ und bestimmen $\ker(f)$ durch
    die Basen $B_1$ und $B_2$.

    Nach Definition $B_1$ gilt für alle $v_1 \in B_1$ $\beta(v_1, v_1) = 0$
    und demnach $f(v_1) = 0$. Es gilt also $v_1 \in \ker(f)$.

    Nach Definition $B_2$ gilt für alle $v_2 \in B_2$ $\beta(v_2, v_2) < 0$
    und demnach $f(v_2) \neq 0$ und $v_2 \notin \ker(f)$.

    Da $\langle B_1 \rangle \cup \langle B_2 \rangle = V$, $f(\langle B_1 \rangle) = 0$
    und $f(\langle B_2 \rangle) \neq 0$, folgt $\langle B_2 \rangle = \mathrm{Bild}(f)$,
    also $\dim(\langle B_2 \rangle) = \mathrm{Rg}(C)$.

    Da ebenfalls die Diagonalmatrix $M_B(\beta)$ den Rang $\dim(\langle B_2 \rangle)$ hat, folgt
    \[
      \mathrm{Rg}(M_B(\beta)) = \mathrm{Rg}(C).
    \]

    Also
    \begin{align*}
      (0,\, \mathrm{Rg}(M_B(\beta)),\, n - \mathrm{Rg}(M_B(\beta)))
        &= (0,\, \mathrm{Rg}(C),\, n - \mathrm{Rg}(C)).
    \end{align*}

    $(0,\, \mathrm{Rg}(C),\, n - \mathrm{Rg}(C))$ ist der Typ von $\beta$.

\end{enumerate}

\end{document}
